\slajdid{09_1-s036}
\begin{frame}[label=09_1-s036]
\gusttekst
Ideja da se umesto otpornika $R_L$ uvede tranzistor. Za početak nMOS tranzistor sa indukovanim kanalom.
\begin{center}\begin{tikzpicture}[x=.8cm,y=.65cm,font=\sitneoznake,>=Latex]
\foreach\y/\lab/\inp in{1.4/L/C,0/D/I}{\begin{scope}[shift={(0,\y)}]
\draw(-.9,0)--(-.28,0);\node[left]at(-.9,0){$V_\inp$};\ifdim\y pt=0pt\draw[fill=white](-.9,0)circle(.05);\fi
\draw(-.28,-.3)--(-.28,.3);\draw(-.2,-.35)--(-.2,.35);\draw(-.2,.25)--(0,.25)--(0,.7);\draw(-.2,-.25)--(0,-.25)--(0,-.7);\node[right]at(.05,0){$T_\lab$};\end{scope}}
\draw(-.2,2.1)--(.2,2.1)node[above left]{$V_{DD}$};\draw(0,-.7)node[ground]{};\draw(0,.7)--(.85,.7)node[ocirc]{}node[right]{$V_O$};
\draw[->,DEplava](2.1,2.6)--(.55,1.45);\end{tikzpicture}
\end{center}
Pošto je bitan režim rada tranzistora $T_L$ (L - load) razmotrićemo dva slučaja u zavisnosti koliki je kontrolni napon $V_C$ na gejtu tranzistora $T_L$. Nadam se da je jasno da je drejn tranzistora $T_L$ na $V_{DD}$ a sors na izlaznom naponu. Isto tako da bi kvalitativno videli šta se dešava smatraćemo da su tranzistori \alert{sa dugačkim kanalom}. Mogli bi da izvedemo uzimajući u obzir i kratki kanal, ali bi ionako u „krajnjim“ režimima uradili uprošćavanja i opet došli do izraza za dugačak kanal, kao što smo videli.
\end{frame}
